\documentclass[a4paper,12pt]{article}

\usepackage[T1]{fontenc}
\usepackage[utf8]{inputenc}
\usepackage[polish]{babel}
\usepackage{geometry}
\usepackage[hidelinks]{hyperref}
\usepackage{listings}
\usepackage{xcolor}
\usepackage{tabularx}
\usepackage{booktabs}
\usepackage{longtable}
\usepackage{array}
\usepackage{parskip}

\geometry{
  a4paper,
  top=2.5cm,
  bottom=2.5cm,
  left=3cm,
  right=2.5cm
}

\lstdefinestyle{bash}{
  basicstyle=\ttfamily\footnotesize,
  breaklines=true,
  frame=single,
  backgroundcolor=\color{gray!8},
  xleftmargin=0.5em,
  xrightmargin=0.5em,
}

\begin{document}

\begin{titlepage}
  \centering

  Collegium Witelona Uczelnia Państwowa\\
  Wydział Nauk Technicznych i Ekonomicznych\\
  Kierunek: Informatyka, semestr VII

  \vspace{4cm}

  {\Large\bfseries Sprawozdanie z pracy projektowej}\\[0.5em]
  z projektu --- aplikacja turystyczna \textbf{,,RTH''}

  \vspace{1cm}

  Przedmiot: Projekt zespołowy

  \vspace{4cm}

  \begin{flushright}
    \textbf{Wykonali:}\\[0.5em]
    Mykyta Klymenko\\
    nr indeksu 44496\\[0.5em]
    \textit{[imię i nazwisko]}\\
    nr indeksu \textit{[numer]}
  \end{flushright}

  \vfill

  Legnica 2027
\end{titlepage}

\tableofcontents
\newpage

\section{Opis przedmiotu zamówienia}

Opis dostarczony przez prowadzącego.

\subsection{Cel projektu}

Przedmiotem zamówienia jest stworzenie aplikacji turystycznej, która umożliwi użytkownikom
dokumentowanie odbytych podróży, prezentowanie ich na interaktywnej mapie oraz przeglądanie
statystyk podróżniczych.

Chcielibyśmy, aby aplikacja pełniła funkcję osobistego dziennika podróży, w którym użytkownik może
zobaczyć, gdzie był, jakimi środkami transportu podróżował i jakie kraje odwiedził.

\subsection{Wymagane funkcjonalności}

\subsubsection*{Konta użytkowników}

\begin{itemize}
  \item Rejestracja, logowanie i wylogowywanie.
  \item Możliwość edycji podstawowych danych profilu.
  \item Ustawienie profilu jako publiczny lub prywatny.
  \item Możliwość przeglądania podróży innych użytkowników, jeśli ich profile są publiczne.
\end{itemize}

\subsubsection*{Zarządzanie podróżami}

\begin{itemize}
  \item Dodawanie, edycja i usuwanie podróży.
  \item Każda podróż powinna mieć nazwę, datę lub zakres dat oraz opcjonalny opis.
  \item Podróż może składać się z dowolnej liczby punktów i odcinków, pokonywanych różnymi
        środkami transportu.
  \item Użytkownik powinien mieć możliwość określenia kolejności odwiedzanych miejsc i wyboru
        środka transportu dla każdego odcinka.
\end{itemize}

Przykładowa podróż:

\begin{enumerate}
  \item Legnica $\rightarrow$ Berlin (samochód)
  \item Berlin $\rightarrow$ Majorka (samolot)
  \item Zwiedzanie Majorki (kilka punktów) (samochód)
  \item Majorka $\rightarrow$ Berlin (samolot)
  \item Berlin $\rightarrow$ Legnica (samochód)
\end{enumerate}

\subsubsection*{Obsługiwane środki transportu}

\begin{itemize}
  \item Samochód --- trasa wyznaczana po rzeczywistej sieci dróg.
  \item Samolot --- połączenie przedstawione linią prostą (choć na mapie powinno przedstawić się
        jako łuk).
  \item Pociąg --- połączenie przedstawione linią prostą.
  \item Inne --- uproszczone połączenie między punktami.
\end{itemize}

\subsubsection*{Interaktywna mapa}

\begin{itemize}
  \item Wyświetlanie pojedynczej podróży na mapie wraz z punktami i odcinkami.
  \item Rozróżnianie środków transportu za pomocą kolorów, ikon lub stylów linii.
  \item Możliwość wyświetlenia wszystkich podróży użytkownika na jednej mapie.
  \item Automatyczne rozpoznawanie i oznaczanie odwiedzonych państw na podstawie zapisanych
        podróży.
  \item Możliwość przybliżania mapy i przeglądania szczegółów tras.
\end{itemize}

\subsubsection*{Statystyki podróżnicze}

Aplikacja powinna generować statystyki zarówno dla całej historii podróży użytkownika, jak i dla
wybranego roku kalendarzowego. Interesują nas przede wszystkim:

\begin{itemize}
  \item Liczba odbytych podróży.
  \item Liczba odwiedzonych państw.
  \item Lista odwiedzonych państw oraz ich wyróżnienie na mapie.
  \item Łączna liczba przebytych kilometrów.
  \item Dystans pokonany poszczególnymi środkami transportu.
  \item Liczba odwiedzonych miejsc.
\end{itemize}

\subsection{Oczekiwany rezultat}

Oczekujemy działającej aplikacji z czytelnym, intuicyjnym interfejsem, pozwalającej użytkownikowi w
prosty sposób dodawać podróże i analizować swoją historię podróżniczą.

Zależy nam przede wszystkim na poprawnym działaniu podstawowych funkcji, sensownej prezentacji tras
na mapie oraz estetycznej i wygodnej obsłudze aplikacji.

\subsection{Uwagi dodatkowe}

Oczekujemy wykorzystania zewnętrznych API i gotowych bibliotek do obsługi map, geokodowania oraz
wyznaczania tras.

Nie oczekujemy systemu nawigacji GPS ani idealnego odwzorowania rzeczywiście przebytej drogi. Trasy
powinny być jednak przedstawiane w sposób wiarygodny i zgodny z wybranym środkiem transportu.

Jesteśmy również otwarci na dodatkowe funkcjonalności zaproponowane przez wykonawców, o ile nie
odbędzie się to kosztem realizacji podstawowych wymagań.

Najważniejsze z perspektywy zamawiającego: chcemy móc po kilku latach korzystania z aplikacji
otworzyć mapę i zobaczyć wszystkie swoje podróże, odwiedzone kraje oraz ciekawe statystyki
pokazujące, jak podróżowaliśmy na przestrzeni lat.

\newpage

\section{Opis technologiczny zaproponowanego rozwiązania}

\subsection{Architektura}

System składa się z trzech części, które komunikują się przez jedno REST API:

\begin{itemize}
  \item \textbf{API} (\texttt{backend/}) --- jedyne miejsce logiki biznesowej i dostępu do bazy danych;
  \item \textbf{panel administracyjny} (\texttt{web/}) --- aplikacja SPA do zarządzania danymi;
  \item \textbf{aplikacja mobilna} (\texttt{mobile/}) --- aplikacja dla użytkowników końcowych.
\end{itemize}

Wszystkie części znajdują się w jednym repozytorium. Środowisko uruchomieniowe opisuje plik
\texttt{docker-compose.yml}.

\subsection{Stos technologiczny}

\begin{tabularx}{\textwidth}{@{}lX@{}}
  \toprule
  \textbf{Warstwa} & \textbf{Technologia} \\
  \midrule
  API & PHP 8.6, Laravel 13, Laravel Sanctum \\
  Dokumentacja API & OpenAPI generowane przez Scramble \\
  Panel administracyjny & Vue 3, TypeScript, Vite, Pinia, Vue Router, Tailwind CSS 4, shadcn-vue \\
  Aplikacja mobilna & \textit{[do uzupełnienia]} \\
  Baza danych & PostgreSQL 19 \\
  Serwer HTTP & nginx 1.31 \\
  Środowisko & Docker Compose, Node.js 26 \\
  CI & GitHub Actions \\
  \bottomrule
\end{tabularx}

\subsection{Kontrola jakości}

Każdy pull request przechodzi w GitHub Actions te same kontrole, które zespół uruchamia lokalnie:

\begin{itemize}
  \item API: formatowanie kodu (Pint), analiza statyczna (Larastan), testy (PHPUnit na PostgreSQL);
  \item panel: lintery (oxlint, ESLint), kontrola typów (vue-tsc), testy (Vitest), build produkcyjny;
  \item tytuł pull requesta zawiera numer zadania z GitHub Issues.
\end{itemize}

\newpage

\section{Podział obowiązków i odpowiedzialności w zespole}

\begin{tabularx}{\textwidth}{@{}lX@{}}
  \toprule
  \textbf{Osoba} & \textbf{Zakres odpowiedzialności} \\
  \midrule
  Mykyta Klymenko & API, panel administracyjny \\
  \textit{[imię i nazwisko]} & panel administracyjny, aplikacja mobilna \\
  \bottomrule
\end{tabularx}

\newpage

\section{Zadania i ich czasochłonność}

Zadania są prowadzone w GitHub Issues. Tabela zawiera czas szacowany przy planowaniu i czas
rzeczywisty po zamknięciu zadania.

\begin{longtable}{@{}p{1cm}p{6.4cm}p{3cm}rr@{}}
  \toprule
  \textbf{Nr} & \textbf{Zadanie} & \textbf{Osoba} & \textbf{Plan [h]} & \textbf{Fakt [h]} \\
  \midrule
  \endhead
  \textit{[\#]} & \textit{[do uzupełnienia]} & & & \\
  \bottomrule
\end{longtable}

\newpage

\section{Instrukcja uruchomienia systemu}

\subsection{Uruchomienie lokalne}

Wymagany jest Docker z Docker Compose.

\begin{lstlisting}[style=bash]
cp backend/.env.example backend/.env
docker compose up -d --build
docker compose exec php composer setup
\end{lstlisting}

Kontener \texttt{web} sam instaluje zależności przy starcie.

\begin{tabularx}{\textwidth}{@{}lX@{}}
  \toprule
  \textbf{Serwis} & \textbf{Adres} \\
  \midrule
  API & \url{http://localhost:8000/api} \\
  Dokumentacja OpenAPI & \url{http://localhost:8000/docs/api} \\
  Panel administracyjny & \url{http://localhost:5173} \\
  Mailpit & \url{http://localhost:8025} \\
  \bottomrule
\end{tabularx}

\subsection{Uruchomienie zdalne}

\textit{[do uzupełnienia po wdrożeniu]}

\end{document}
